\section{Sorting}
\textbf{\textcolor{Blue}{Input}}: An array $A$ of $n$ elements, and a total order $\le$ on the elements.\\
\textbf{\textcolor{Blue}{Output}}: A permutation of $A$ such that $A[1]\le A[2]\le \ldots \le A[n]$.

\paragraph{Properties of Sorting Algorithms}
\textcolor{Blue}{Time complexity}: in the worst case, average case, or best
case.

\textcolor{Blue}{Stability}: Original ordering is preserved for elements of
equal values.

\textcolor{Blue}{In-place}: Very little extra space is used (e.g. $O(1)$ or
$O(\log n)$).

\subsection{Comparison-based Sorting}
\textcolor{Blue}{Comparison-based sorting}: Sorting algorithms that
only uses comparisons (and no other information) to determine the order of
elements.

\paragraph{Lower bound for Comparison-based Sorting}
\textbf{Claim}: The \textcolor{Blue}{worst-case time complexity} of any
comparison-based sorting algorithm is $\Omega(n\log n)$.

\textbf{Proof}:
Comparison-based sorting can be modelled as a
\textbf{\textcolor{Bittersweet}{binary decision tree}}. Each internal node
represents a \textcolor{Bittersweet}{comparison} between two elements, and each
leaf represents a possible permutation of the input.

The worst-case runing time $\ge $ worst-case number of comparisons = height of
the tree.

Since the \textcolor{Bittersweet}{binary tree has at least $n\mathpunct{!}$
leaves}, we have the height of the tree $h\ge \log{\left( n\mathpunct{!}
\right) }$. Using Stirling's approximation, we have $\log{\left( n\mathpunct{!}
\right) } = \Theta(n\log n)$.

\subsection{Average-case Analysis for Quick Sort}
If worst-case time complexities do not capture an algorithm's typical
performance, we use \textcolor{Blue}{average-case analysis} for it's expected
time complexity.

\textbf{Claim}: The average-case time complexity of Quick Sort is $O(n\log n)$.

\textbf{Proof}: WLOG, assume all elements are distinct. Suppose
$a_1<a_2<\cdots<a_{n}$. The input array is a \textcolor{Bittersweet}{random
permutation $\pi$} of $\left( a_1,a_2,\ldots,a_{n} \right) $. Suppose $\pi$ is
\textcolor{Bittersweet}{\textbf{uniformly random}}. That is, each of the
$n\mathpunct{!}$ permutations is equally likely to be the input.

Then the pivot is equally likely to be any of the $n$ elements, and the
permutations for each of the recursive calls are also uniformly random. Then
the \textcolor{Blue}{expected running time $A(n)$} is given by the recurrence:
\begin{align*}
  A(n)&=\frac{1}{n}\sum_{j=1}^{n} \left[ A\left( j-1 \right) +A\left( n-j
\right)  \right] +cn\\
      &=\frac{2}{n}\sum_{j=0}^{n-1} A(j) +cn
\end{align*}
Solving the recurrence, we have $A(n) \in O(n log n)$.
